\chapter{Grundlagen}
\label{ch:grundlagen}

Im folgenden Abschnitt werden grundlegende Begriffe f{\"u}r das Verst{\"a}ndnis dieser Arbeit definiert. Dabei werden unter anderem die Begriffe \emph{Job} und \emph{Task} verwendet. Ein Job ist eine bestimmte Ausf{\"u}hrung einer Aufgabe. Ein Task beschreibt auf der anderen Seite die dazugeh{\"o}rige Folge oder Klasse von Jobs. Bei einem periodischen Task entstehen solche Jobs in bestimmten Zeitabst{\"a}nden (Liu, 2000, S. 27--28). Bspw. k{\"o}nnte ein Task wie folgt aussehen: \glqq{}Lese alle 10 ms einen Sensor aus.\grqq{}, wobei die einzelnen Jobs unterteilt werden:

\begin{itemize}
\item Job 1 bei t = 0 ms
\item Job 2 bei t = 10 ms
\item Job 3 bei t = 20 ms
\item (\ldots{})
\end{itemize}
\section{Systeme und zeitliche Anforderungen}
\textbf{System} Ein System l{\"a}sst sich allgemein {\"u}ber den Zusammenhang zwischen einer Menge von Eingaben und einer daraus entstehenden Menge an Ausgaben beschreiben. Wie dabei dieser Zusammenhang intern entsteht, ist dabei oft nebens{\"a}chlich. Ein Thermostat nimmt bspw. als Eingabe die gemessene Raumtemperatur und macht daraus ein heiz- oder K{\"u}hlsignal als Ausgabe. Dabei ist von au{\ss}en komplett unbekannt, wie die interne Regellogik aufgebaut ist. Dabei wird solch ein Fall Black-Box-Betrachtung genannt. (Laplante \& Ovaska, 2012, S.3)

\textbf{Ereignis }Laplante und Ovaska benutzen den Begriff Ereignis (\emph{event}) f{\"u}r einen Wechsel im Programmablauf, bei dem die Ausf{\"u}hrung nicht einfach mit der n{\"a}chsten sequenziellen Instruktion fortgesetzt wird. Solche Wechsel k{\"o}nnen dabei aus dem Programm selbst entstehen oder von au{\ss}en ausgel{\"o}st werden. Zudem unterscheiden die Autoren synchrone und asynchrone Ereignisse. Eine programmgesteuerte Verzweigung geh{\"o}rt zur ersten Gruppe, eine von au{\ss}en eintreffende Unterbrechung zur zweiten. Die Freigabe eines Jobs kann zudem mit einem solchen Ereignis verbunden sein (Laplante \& Ovaska, 2012, S. 10).

\textbf{Freigabezeitpunkt }Der Freigabezeitpunkt (\emph{release} \emph{time}) gibt an, ab wann ein Job zur Ausf{\"u}hrung bereitsteht. Er ist jedoch vom echten Beginn der Verarbeitung zu unterscheiden, da der Job nach seiner Freigabe noch auf Rechenzeit warten kann. Die Freigabe kann bspw. durch das Eintreffen von neuen Sensordaten ausgel{\"o}st werden. Ein Ereignis und ein Job sind dabei nicht gleich. Das Ereignis ist der ausl{\"o}sende Vorgang, w{\"a}hrend der Job die anschlie{\ss}end zu bearbeitende Arbeitseinheit ist (Laplante \& Ovaska, 2012, S. 10; Liu, 2000, S. 27).

\textbf{Abschlusszeitpunkt }Der Abschlusszeitpunkt (\emph{completion time}) eines Jobs bezeichnet den Zeitpunkt, an welchem dieser fertig umgesetzt wurde (Liu, 2000, S.28).

\textbf{Antwortzeit} Wenn ein System auf eine Eingabe reagiert, dann vergeht eine gewisse Zeit, bis das System das angeforderte Verhalten und alle dazugeh{\"o}rigen Ausgaben komplett realisiert hat. Genau dieser Zeitraum zwischen dem Anlegen der Eingabe und der kompletten Reaktion wird Antwortzeit (\emph{response} \emph{time}) genannt. Bei dem Thermostat-Beispiel w{\"a}re das bspw. die Zeit zwischen dem Erfassen einer niedrigen Raumtemperatur und dem Einschalten der Heizung. Welche Antwortzeit dabei akzeptable ist, h{\"a}ngt vom Anwendungsfall ab. Eine Sicherheitsabschaltung darf z.B. weniger Zeit in Anspruch nehmen als eine Komfortfunktion (Laplante \& Ovaska, 2012, S.4). F{\"u}r die Antwortzeit eines einzelnen Jobs werden dessen Freigabe und Abschluss als zeitliche Bezugspunkte verwendet. Wenn \(r_i\) den Freigabezeitpunkt bezeichnet und \(f_i\) den Abschlusszeitpunkt, ergibt sich die Antwortzeit als \(R_i = f_i - r_i\). Sie beinhaltet damit sowohl m{\"o}gliche Wartezeiten als auch die eigentliche Ausf{\"u}hrung des Jobs (Liu, 2000, S. 28).

\textbf{Frist }Eine Frist (\emph{deadline}) legt fest, bis zu welchem Zeitpunkt ein Job abgeschlossen sein muss. Die relative Frist \(D_i\) beschreibt die verf{\"u}gbare Zeitspanne ab seiner Freigabe. Zusammen mit dem Freigabezeitpunkt \(r_i\) ergibt sich daraus die absolute Frist \(d_i = r_i + D_i\). Ob die zeitliche Vorgabe eingehalten wurde, l{\"a}sst sich somit durch den Vergleich mit dem Abschlusszeitpunkt \(f_i\) bestimmen: Der Job ist dabei fristgerecht beendet, wenn \(f_i \leq d_i\) gilt (Liu, 2000, S.27-28).

\textbf{Frist{\"u}berschreitung }Die Frist{\"u}berschreitung (\emph{tardiness}) gibt an, wie lange ein Job nach seiner absoluten Frist noch bis zum Abschluss ben{\"o}tigt. Sie wird aus Abschlusszeitpunkt \(f_i\) und absoluter Frist \(d_i\) als \(T_i = \max(0, f_i - d_i)\) berechnet. Bei einem rechtzeitig beendeten Job ist \(T_i = 0\). Nur ein versp{\"a}teter Abschluss f{\"u}hrt zu einem positiven Wert (Liu, 2000, S. 28).

\textbf{Zeitliche Anforderung} Zeitliche Anforderungen (\emph{timing} \emph{constraints}) legen fest, welche Zeitpunkte oder Zeitr{\"a}ume ein Job bei seiner Ausf{\"u}hrung einhalten muss. Eine einfache Form ergibt sich dabei aus Freigabezeitpunkt und Frist. Anwendungen k{\"o}nnen jedoch weitere Bedingungen vorgeben. F{\"u}r ihre Bewertung ist es relevant, welche Folgen eine Frist{\"u}berschreitung hat. Bei einer weichen zeitlichen Anforderung bleibt ein versp{\"a}tetes Ergebnis an sich verwendbar, verliert aber an Nutzen. Eine verz{\"o}gert aktualisierte Sensoranzeige kann bspw. weiterhin Informationen liefern. Bei einer harten zeitlichen Anforderung kann ein versp{\"a}tetes Ergebnis jedoch unbrauchbar oder sogar sch{\"a}dlich sein, bspw. wenn eine Schutzabschaltung zu sp{\"a}t erfolgt (Liu, 2000, S. 28).

\textbf{Ausgefallenes System} Ein System gilt genau dann als ausgefallen (\emph{failed} \emph{system}), wenn es eine oder mehrere Anforderungen nicht mehr erf{\"u}llt, welche in seiner Spezifikation festgelegt wurden. Dies gilt dabei unabh{\"a}ngig davon, ob die Ursache in der Hardware, der Software oder im zeitlichen Verhalten liegt. (Laplante \& Ovaska, 2012, S.5)

\textbf{Deterministisches System }Ein System hei{\ss}t deterministisch (\emph{deterministic} \emph{system}), wenn sein Verhalten komplett vorhersagbar ist. Wenn sich das System hierbei in einem bestimmten Zustand befindet und es eine bestimmte Eingabe erh{\"a}lt, dann steht fest, welche Ausgabe es liefert und in welchen Zustand es als N{\"a}chstes wechselt. Dabei gibt es in diesem Sinne keinen Zufall und keine Mehrdeutigkeit im Verhalten. Eine schw{\"a}chere, aber verwandte Eigenschaft ist der Ereignisdeterminismus. Hierbei muss nur f{\"u}r die Eingaben, die ein Ereignis ausl{\"o}sen, feststehen, wie das System reagiert. Jedes deterministische System ist also dadurch automatisch auch ereignisdeterministisch, andersherum jedoch nicht. (Laplante \& Ovaska, 2012, S.11)

\section{Real Time Systems}
\textbf{Echtzeitsystem} Bei einem Echtzeitsystem (\emph{real-time system, RTS}) reicht ein Ergebnis, welches inhaltlich richtig ist, nicht allein aus. Es muss au{\ss}erdem innerhalb der Zeit verf{\"u}gbar sein, die f{\"u}r die Anwendung vorgegeben ist. Die zeitliche Anforderung geh{\"o}rt daher zur Korrektheitsbedingung des Systems. Welche Folgen dabei eine versp{\"a}tete Reaktion hat, h{\"a}ngt von der Anwendung und ihren Anforderungen ab (Laplante \& Ovaska, 2012, S. 5).

Echtzeitsysteme stehen oft in einer laufenden Wechselwirkung mit ihrer Umgebung. Eingaben k{\"o}nnen neue Verarbeitungsaufgaben ausl{\"o}sen, deren Ergebnisse auf der anderen Seite k{\"o}nnen auf die Umgebung einwirken. Diese reaktive Arbeitsweise sorgt daf{\"u}r, dass periodisch geplante Aufgaben eingeschlossen werden (Laplante \& Ovaska, 2012, S. 5).

Bei einem Temperaturregler k{\"o}nnte bspw. festgelegt werden, dass innerhalb einer vorgegebenen Frist reagiert werden muss, wenn eine Grenzwert{\"u}berschreitung erkannt wird. Ob die Temperatur daf{\"u}r regelm{\"a}{\ss}ig oder ereignisgesteuert erfasst wird, ist hierbei eine getrennte Entwurfsentscheidung.

\textbf{Weiches Echtzeitsystem} Bei weicher Echtzeit beeintr{\"a}chtigt eine versp{\"a}tete Reaktion die Leistungs- oder Dienstqualit{\"a}t, ohne die Funktion des Systems direkt aufzuheben (Laplante \& Ovaska, 2012, S. 6).

\textbf{Hartes Echtzeitsystem} Bei harter Echtzeit kann schon eine einzelne verfehlte Frist zu einem kompletten oder katastrophalen Ausfall f{\"u}hren (Laplante \& Ovaska, 2012, S. 6).

\textbf{Festes Echtzeitsystem} Bei fester Echtzeit k{\"o}nnen wenige Fristverletzungen toleriert werden. Eine gr{\"o}{\ss}ere Zahl kann dabei zum Systemausfall f{\"u}hren (Laplante \& Ovaska, 2012, S. 7).

\section{Eingebettete Systeme}
\textbf{Eingebettetes System} Bei einem eingebetteten System werden Hardware und Software auf eine bestimmte Aufgabe oder eine bestimmte Menge an Funktionen begrenzt. Dies stellt hierbei auch den Unterschied des Entwurfsziels von dem eines PCs dar, der f{\"u}r mehrere unterschiedliche Anwendungen benutzt werden kann. Das unterscheidet dabei zwischen einem zweckgebundenen Rechner (\emph{special-purpose computer}) und einem Universalrechner (\emph{general-purpose computer}). Die zum Ger{\"a}t dazugeh{\"o}rende Software wird zudem Firmware genannt (Das, 2013, S. 34).

F{\"u}r die Beschreibung des Aufbaus verwendet Das ein simples Modell aus Prozessor, Speicher, Sensoren und Aktoren. Sensoren stellen Eingabedaten bereit, der Prozessor verarbeitet sie und Aktoren machen eine Reaktion auf die Umgebung m{\"o}glich. Der Speicher h{\"a}lt die Programme und Daten, die f{\"u}r den Betrieb gebraucht werden (Das, 2013, S. 37--38).

Ein Beispiel w{\"a}re ein Temperaturregler. Ein Sensor liefert einen Temperaturwert. Das Programm vergleicht ihn dabei mit einem vorgegebenen Sollwert und entscheidet dann anschlie{\ss}end, ob eine Heizung ein- oder ausgeschaltet werden soll.

Die zeitlichen Anforderungen eines eingebetteten Systems basieren auf der Wechselwirkung mit der physischen Umgebung. Sensoren bestimmen dabei, wann neue Zustandsinformationen verf{\"u}gbar werden. Aktoren m{\"u}ssen innerhalb der f{\"u}r den jeweiligen Prozess vorgegebenen Zeit reagieren. Aus diesen Eigenschaften werden die Fristen der dazugeh{\"o}rigen Verarbeitungsaufgaben festgelegt. Ob dabei harte oder weniger strenge Echtzeitanforderungen gebraucht werden, h{\"a}ngt von den m{\"o}glichen Folgen einer versp{\"a}teten Reaktion ab (Liu, 2000, S. 30).

\textbf{Eingebettetes Betriebssystem }Ein eingebettetes System braucht nicht f{\"u}r jede Anwendung ein Betriebssystem. Bei einer einfachen Softwareorganisation kann eine sich wiederholende durchlaufene Hauptschleife die Aufgaben bearbeiten. Ein solcher Superloop kann zudem durch Interrupt-Routinen erweitert werden, die auf bestimmte Ereignisse reagieren. Die Verwendung von Interrupts allein macht dabei noch keinen Superloop aus. Der entscheidende Punkt ist hierbei die Organisation der Verarbeitung in einer wiederkehrenden Schleife (Das, 2013, S. 296).

Wenn jedoch mehrere Aufgaben und gemeinsam genutzte Ressourcen koordiniert werden m{\"u}ssen, dann kann ein Betriebssystem diese Verwaltung {\"u}bernehmen. Es stellt dann grundlegende Dienste bereit, auf denen die Anwendungssoftware aufbaut. Ein Betriebssystem, das f{\"u}r solche Einsatzbedingungen benutzt wird, wird eingebettetes Betriebssystem (\emph{embedded operating system}) genannt (Das, 2013, S. 41, 296).

\textbf{Mikrocontroller }Ein Mikrocontroller (Microcontroller Unit, MCU) integriert eine Recheneinheit und weitere f{\"u}r Komponenten, die f{\"u}r den Ger{\"a}tebetrieb gebraucht werden, auf einem Chip. Dazu geh{\"o}ren Speicher, Zeitgeber und bspw. serielle und parallele Schnittstellen. Dadurch lassen sich viele Aufgaben ohne einen Aufbau aus getrenntem Prozessor, externem Speicher und zus{\"a}tzlichen Schnittstellenbausteinen realisieren. Nichtfl{\"u}chtiger Speicher kann den Programmcode halten, w{\"a}hrend RAM f{\"u}r sich ver{\"a}ndernde Daten w{\"a}hrend der Ausf{\"u}hrung gebraucht wird (Das, 2013, S. 39).

Das beschreibt dar{\"u}ber hinaus auch die Integration von weiteren Systemfunktionen auf einem Chip im Kontext zu dem Begriff \glqq{}System on Chip\grqq{} (SoC). F{\"u}r diese Arbeit ist vor allem die gemeinsame Bereitstellung von Prozessor, Speicher und Peripheriefunktionen auf dem verwendeten Mikrocontroller relevant (Das, 2013, S. 39).

\section{Echtzeitbetriebssysteme}
\textbf{Betriebssystem }Ein Betriebssystem (\emph{operating system}, \emph{OS}) {\"u}bernimmt grundlegende Verwaltungsaufgaben, die sonst von den einzelnen Anwendungen selbst gel{\"o}st werden m{\"u}ssten. Dazu geh{\"o}ren die Nutzung der CPU, die Speicherverwaltung und der Zugriff auf Ein- und Ausgabeger{\"a}te. Dabei stellt es daf{\"u}r Dienste und Schnittstellen bereit, damit Anwendungen nicht die Einzelheiten der benutzten Hardware ber{\"u}cksichtigen m{\"u}ssen. Diese Abstraktion macht die Nutzung der Ressourcen leichter und trennt dabei anwendungsspezifische Aufgaben von grundlegenden Systemfunktionen. Das Betriebssystem geh{\"o}rt damit zur Systemsoftware, wobei Anwendungsprogramme auf deren Diensten aufbauen (Das, 2013, S. 222).

\textbf{Kernel} Der Kernel stellt die grundlegenden Laufzeitdienste eines Betriebssystems bereit. Dazu z{\"a}hlen bspw. Aufgabenplanung, Interrupt-Behandlung und Speicherverwaltung. Das zeigt hierbei seine Stellung durch ein Zwiebelmodell. Um den Kernel herum liegen weitere Teile der Systemsoftware, Entwicklungswerkzeuge und anschlie{\ss}end die Anwendungsprogramme. Die Darstellung ordnet zudem diese Bausteine nach ihrer Rolle im Gesamtsystem (Das, 2013, S. 231--232).

F{\"u}r Betriebssysteme, die eine Trennung zwischen Kernel- und Nutzerbereich implementieren, beschreibt Das au{\ss}erdem verschiedene Zugriffsrechte. Privilegierte Operationen werden vom Kernel ausgef{\"u}hrt. Anwendungen k{\"o}nnen seine Dienste {\"u}ber Systemaufrufe (System Calls) anfordern. Dadurch l{\"a}sst sich der Zugriff auf grundlegende Systemressourcen kontrollieren. Ob und wie diese Trennung dabei umgesetzt wird, h{\"a}ngt von der Betriebssystem- und Hardwarekonfiguration ab (Das, 2013, S. 232).

\textbf{Kernel-Design }Monolithische und mikrokernelbasierte Entw{\"u}rfe unterscheiden sich vor allem darin, wie sie Betriebssystemdienste auf privilegierte und getrennt ausgef{\"u}hrte Komponenten verteilen. Bei einem monolithischen Kernel sind dabei einige Dienste im Kernelbereich zusammengefasst. Ein Fehler in einer dieser Komponenten kann dadurch auch andere Teile des Systems beeinflussen. Das nennt Linux als Beispiel f{\"u}r diese Architekturvariante (Das, 2013, S. 233--234).

Ein Mikrokernel schr{\"a}nkt den Kern auf der anderen Seite auf grundlegende Funktionen ein. Zudem verlagert es weitere Dienste in getrennte Komponenten bzw. Serverprozesse. Diese Trennung kann die Ausbreitung von Fehlern eingrenzen, braucht jedoch eine Kommunikation zwischen den beteiligten Diensten. Daraus entsteht eine Balancesuche zwischen Fehlerisolation, Modularit{\"a}t und Kommunikationsaufwand (Das, 2013, S. 234).

Die f{\"u}r die Untersuchung relevanten Eigenschaften des Zephyr-Kernels werden im folgenden Abschnitt anhand seiner Dokumentation beschrieben.

\textbf{Echtzeitbetriebssystem} Ein Echtzeitbetriebssystem (\emph{real-time operating system, RTOS}) stellt wie andere Betriebssysteme Dienste zur Verwaltung von Aufgaben und Hardware-Ressourcen bereit. F{\"u}r seinen Einsatz sind jedoch noch die zeitlichen Anforderungen der Anwendung relevant. Die Aufgabenplanung muss deshalb darauf achten, welche Verarbeitungen rechtzeitig abgeschlossen werden sollen. Das hebt Scheduling als einen relevanten Unterschied im Vergleich zu allgemeinen Betriebssystemen hervor (Das, 2013, S. 296--297).

F{\"u}r diese Arbeit ist daran vor allem relevant, wie die Priorisierung von Aufgaben das zeitliche Verhalten des Systems beeinflusst. Die genaue Auswahlregel des Zephyr-Schedulers wird in Abschnitt 2.5.1a erl{\"a}utert.

\section{Zephyr RTOS}
Zephyr ist ein quelloffenes (open-source) RTOS, das f{\"u}r den Einsatz auf ressourcenbeschr{\"a}nkten eingebetteten Systemen entwickelt wurde. Das Projekt wird von der Linux Foundation als ein gemeinschaftlich entwickeltes Open-Source-Projekt unterst{\"u}tzt. Zephyr verfolgt das Ziel, ein skalierbares und plattform{\"u}bergreifendes OS f{\"u}r unterschiedliche Embedded-Anwendungen bereitzustellen. Zephyr kann, durch die Unterst{\"u}tzung unterschiedlicher Prozessorarchitekturen und Hardwareplattformen auf Systemen mit verschiedenen Ressourcen- und Leistungsanforderungen eingesetzt werden. Die Architektur ist modular und konfigurierbar aufgebaut. Dadurch k{\"o}nnen ben{\"o}tigte Komponenten deaktiviert werden. Der Zephyr-Kernel unterst{\"u}tzt zudem kooperatives und pr{\"a}emptives Threading und stellt Funktionen f{\"u}r die Aufgabenverwaltung, Synchronisation, Speicherverwaltung und Interaktion mit Hardware bereit. Au{\ss}erdem verf{\"u}gt Zephyr {\"u}ber integrierte Schnittstellen f{\"u}r Ger{\"a}tetreiber und Subsysteme f{\"u}r bspw. Metzwerkkommunikation, Bluetooth, USB, Dateisysteme und Logging. Durch die M{\"o}glichkeit, Zephyr zu konfigurieren, kann der ben{\"o}tigte Speicher- und Funktionsumfang an die verf{\"u}gbare Hardware angepasst werden. Auf der anderen Seite unterst{\"u}tzt Zephyr auch leistungsf{\"a}higere Systeme. (Zephyr Project, 2024, \glqq{}About the Zephyr Project\grqq{})

Zephyr basiert auf einem Kernel mit geringem Speicherbedarf. Dieser wurde f{\"u}r ressourcenbeschr{\"a}nkte und eingebettete Systeme entwickelt. Der Zephyr Kernel unterst{\"u}tzt hierbei eine Vielzahl von CPU-Architekturen, darunter die in dieser Arbeit verwendete ARMv7E-M (Cortex-M4F) Architektur. Zephyr wurde zudem durch Subsysteme modular konzipiert. Subsysteme umfassen hierbei logisch abgegrenzte Teile des OS, welche bestimmte Dienste bereitstellen, darunter bspw. Netzwerke, Dateisysteme oder Kommunikationsprotokolle, etc. Jedes Subsystem ist modular aufgebaut und kann somit konfiguriert, angepasst und erweitert werden. (Zephyr Project, 2024, \glqq{}Introduction\grqq{})

Die Messanwendung dieser Arbeit ist in C implementiert und nutzt dabei die {\"u}ber Zephyrs Header bereitgestellten Funktionen und Makros. Welche Dienste verf{\"u}gbar sind, h{\"a}ngt dabei auch von der ausgew{\"a}hlten Konfiguration und Zielplattform ab. F{\"u}r den Anwendungseinstieg schreibt die Dokumentation den R{\"u}ckgabetyp \lstinline[style=fileinline]|int| f{\"u}r \lstinline[style=fileinline]|main()| vor. Beim Beenden soll die Funktion den Wert null zur{\"u}ckgeben. Andere R{\"u}ckgabewerte sind n{\"a}mlich bereits reserviert (Zephyr Project, 2023, \glqq{}C Language Support\grqq{}).

\subsection{Laufzeit-Konzepte}
Der Zephyr-Kernel bildet das Herzst{\"u}ck jeder Zephyr-Anwendung. Er stellt dabei eine ressourcenschonende, performante mehr-Thread-Ausf{\"u}hrungsumgebung mit einem umfangreichen Funktionsangebot bereit. Dabei bauen darauf mehrere weitere Bestandteile des Systems auf, darunter Ger{\"a}tetreiber, Netzwerk-Stacks und anwendungsspezifischer Code. Durch die Konfigurierbarkeit l{\"a}sst sich der Kernel auf genau die Funktionen beschr{\"a}nken, die eine Anwendung wirklich ben{\"o}tigt. Dadurch ist er sowohl f{\"u}r Systeme mit geringem Speicherbedarf als auch f{\"u}r deutlich komplexere Anwendungen mit mehreren Kommunikationsstellen und aufwendigerer Nebenl{\"a}ufigkeit geeignet.

Die Zephyr-Dokumentation unterteilt dabei die Kernel-Dienste in mehrere Kategorien:

\begin{itemize}
\item \glqq{}Scheduling, Interrupts and Synchronization\grqq{}, darunter Threads, Scheduling, Interrupts, Semaphores and Synchronization
\item \glqq{}Data Passing\grqq{}, darunter Mechanismen zum Datenaustausch zwischen Threads und ISRs
\item \glqq{}Memory Management\grqq{}
\item \glqq{}Timing\grqq{}, darunter Zeitgeber und Uhren
\end{itemize}
F{\"u}r die Untersuchung wurden aus diesen Kategorien besonders die folgenden Themen ausgesucht, da sie die Grundlagen f{\"u}r die sp{\"a}ter gemessenen Benchmarks bilden:

\begin{itemize}
\item Threads und Scheduling
\item Synchronisation mittels Semaphoren
\item Interrupts
\item Timer und die zugeh{\"o}rige Zeitbasis des Kernels
\item Speicherverwaltung
\end{itemize}
Erg{\"a}nzend zu diesen Themen wird zudem das Logging-Subsystem definiert, welches als eigenst{\"a}ndiger Dienst au{\ss}erhalb des Kernel-Services in der Zephyr-Dokumentation dokumentiert wird. Au{\ss}erdem wird die Timing-API zur Zeitmessung behandelt, da diese und das Logging-Subsystem beide f{\"u}r die Konfiguration bzw. f{\"u}r die Mess-Infrastruktur dieser Arbeit wichtig sind.

\subsubsection{Threads und Scheduling}
F{\"u}r den Kontextwechsel-Benchmark ist vor allem wichtig, wann Zephyr die Auswahl des laufenden Threads neu pr{\"u}ft. Solche Scheduling-Gelegenheiten entstehen bspw., wenn ein Thread blockiert, ein anderer Thread bereit wird oder der laufende Thread die CPU {\"u}ber \lstinline[style=fileinline]|k_yield()| freiwillig abgibt. Auch die R{\"u}ckkehr aus einer Interrupt-Behandlung kann hierbei zu einer neuen Auswahl f{\"u}hren. Die Dokumentation nennt diese Stellen Reschedule Points (Zephyr Project, 2024, \glqq{}Scheduling\grqq{}, Abschnitt \glqq{}Concepts\grqq{}).

Ob ein anderer Thread ausgef{\"u}hrt wird, h{\"a}ngt dabei von den bereiten Threads, ihren Priorit{\"a}ten und den Scheduling-Regeln ab, die gerade gelten. Bei einem Wechsel m{\"u}ssen die Zustandsinformationen erhalten bleiben, die f{\"u}r die sp{\"a}tere Fortsetzung gebraucht werden. Die in Kapitel 5.1 untersuchten Varianten benutzen daf{\"u}r die Aktivierung eines wartenden Threads bzw. die freiwillige CPU-Abgabe.

\textbf{Auswahlkriterium des} \textbf{Schedulers} Der Scheduler w{\"a}hlt den bereiten Thread mit der h{\"o}chsten Priorit{\"a}t als aktuellen Thread aus. Existieren mehrere bereite Threads derselben Priorit{\"a}t, wird derjenige gew{\"a}hlt, der am l{\"a}ngsten wartet. Die Ausf{\"u}hrung von Interrupt-Service-Routinen (ISRs) hat dabei Vorrang vor der Ausf{\"u}hrung jedes Threads. Dies gilt hierbei sowohl f{\"u}r kooperative als auch pr{\"a}emptive Threads.

\textbf{Kooperative Threads }Kooperative und pr{\"a}emptive Threads unterscheiden sich darin, unter welchen Bedingungen andere Threads ihre Ausf{\"u}hrung verdr{\"a}ngen k{\"o}nnen. Ein kooperativer Thread gibt die CPU im regul{\"a}ren Scheduling selbst ab, bspw. durch Blockieren oder durch \lstinline[style=fileinline]|k_yield()|. Bei einem Yield bekommen andere bereite Threads gleicher oder h{\"o}herer Priorit{\"a}t die M{\"o}glichkeit, ausgef{\"u}hrt zu werden. Wenn kein geeigneter Thread vorhanden ist, dann kann der aufrufende Thread direkt weiterlaufen. L{\"a}ngere Berechnungen ohne freiwillige Abgabe k{\"o}nnen deshalb andere Threads verz{\"o}gern (Zephyr Project, 2024, \glqq{}Scheduling\grqq{}, Abschnitt \glqq{}Cooperative Time Slicing\grqq{}).

\textbf{Pr{\"a}emptive Threads }Ein pr{\"a}emptiver Thread kann auf der anderen Seite verdr{\"a}ngt werden, wenn ein h{\"o}herpriorer Thread breit ist. Dar{\"u}ber hinaus kann Zephyr die CPU-Zeit zwischen gleichprioren pr{\"a}emptiven Threads {\"u}ber Zeitscheiben verteilen. Nach Ablauf einer Zeitscheibe bekommen andere passende Threads dieser Priorit{\"a}t eine M{\"o}glichkeit, ausgef{\"u}hrt zu werden. Ob und f{\"u}r welche Priorit{\"a}ten diese Rotation benutzt wird, ist dabei konfigurierbar (Zephyr Project, 2024, \glqq{}Scheduling\grqq{}, Abschnitt \glqq{}Preemptive Time Slicing\grqq{}).

Diese Unterscheidung ist f{\"u}r die Thread-Auswahl relevant. Sie bedeutet nicht, dass kooperative Threads von Grund auf vor Interrupts gesch{\"u}tzt sind.

\textbf{Scheduler Locking }Ein pr{\"a}emptiver Thread, der w{\"a}hrend einer kritischen Operation nicht unterbrochen werden m{\"o}chte, kann den Scheduler {\"u}ber \lstinline[style=fileinline]|k_sched_lock()| anweisen, ihn wie einen kooperativen Thread zu behandeln. Danach stellt \lstinline[style=fileinline]|k_sched_unlock()| den urspr{\"u}nglichen pr{\"a}emptierbaren Status wieder her.

\textbf{Thread Sleeping und Busy Waiting }Der Kernel stellt, au{\ss}er dem bereits beschriebenen \lstinline[style=fileinline]|k_sleep(),| eine zweite M{\"o}glichkeit durch \lstinline[style=fileinline]|k_busy_wait()| bereit, die Ausf{\"u}hrung f{\"u}r eine bestimmte Zeit zu verz{\"o}gern, ohne die CPU dabei an einen anderen Thread abzugeben. Laut Zephyr-Dokumentation wird ein Busy-Wait normalerweise dann verwendet, wenn die Verz{\"o}gerung zu kurz ist, um den Aufwand eines Kontextwechsels zu einem anderen Thread und zur{\"u}ckzurechtfertigen.

\textbf{Empfohlene Einsatzgebiete} Die Zephyr-Dokumentation empfiehlt kooperative Threads unter anderem f{\"u}r Ger{\"a}tetreiber und andere leistungs-kritische Arbeiten. Diese Threads werden au{\ss}erdem auch zur gegenseitigen Ausschlusssicherung ohne ein eigenes Kernel-Objekt, bspw. einem Mutex, empfohlen. Auf der anderen Seite werden pr{\"a}emptive Threads empfohlen, um zeitkritischer Verarbeitung Vorrang vor weniger zeitkritischer Verarbeitung zu geben.

\subsubsection{Synchronisation (Semaphore)}
Ein Semaphor macht es m{\"o}glich, die Verf{\"u}gbarkeit einer Ressource oder die Freigabe eines Verarbeitungsschritts zwischen Ausf{\"u}hrungskontexten zu organisieren. Zephyr stellt daf{\"u}r einen z{\"a}hlenden Semaphor bereit. Sein Z{\"a}hler beschreibt die aktuell verf{\"u}gbare Anzahl von Freigaben. Ein Grenzwert begrenzt den maximalen Z{\"a}hlerstand. Bei der Initialisierung darf der Z{\"a}hler auch null sein, muss aber innerhalb des zugelassenen Bereichs bis zum Grenzwert liegen (Zephyr Project, 2022, \glqq{}Semaphores\grqq{}, Abschnitt \glqq{}Concepts\grqq{}).

Eine erfolgreiche Entnahme verbraucht dabei eine Freigabe. Wenn keine verf{\"u}gbar ist, dann kann ein Thread auf eine sp{\"a}tere Freigabe warten. Bei mehreren wartenden Threads achtet Zephyr zuerst auf die Priorit{\"a}t und bei gleicher Priorit{\"a}t auf die Wartezeit. Ein Semaphor kann dabei auch aus einer ISR gegeben werden. Eine Entnahme aus einer ISR muss dabei nichtblockierend mit \lstinline[style=fileinline]|K_NO_WAIT| erfolgen (Zephyr Project, 2022, \glqq{}Semaphores\grqq{}, Abschnitte \glqq{}Concepts\grqq{} und \glqq{}API Reference\grqq{}).

Im Kontextwechsel-Benchmark wird das Semaphor zuerst leer angelegt. Der Worker wartet anschlie{\ss}end darauf, dass der Hauptthread eine Freigabe ausl{\"o}st. Das Semaphor dient in diesem Fall dazu, einen definierten {\"U}bergang zu synchronisieren und nicht dazu, mehrere gleichartige Ressourcen zu verwalten.

\textbf{Zwei Einsatzmuster} Laut der Zephyr-Dokumentation kann ein \glqq{}volles\grqq{} Semaphor (Z{\"a}hlerstand gleich Grenzwert) daf{\"u}r verwendet werden, die Anzahl der Threads zu begrenzen, die gleichzeitig einen kritischen Abschnitt ausf{\"u}hren d{\"u}rfen. Ein \glqq{}leeres\grqq{} Semaphor (Z{\"a}hlerstand gleich Null, Grenzwert gr{\"o}{\ss}er Null) erzeugt auf der anderen Seite eine Art Schranke, an der kein wartender Thread vorbei darf, bis das Semaphor erh{\"o}ht wird. Dieses zweite Muster bildet die Grundlage f{\"u}r die Freigabe eines wartenden Threads durch einen Erzeuger-Thread oder eine ISR.

\textbf{Definition und Nutzung im Code }Ein Semaphor wird {\"u}ber eine Variable vom Typ \lstinline[style=fileinline]|k_sem| angelegt und anschlie{\ss}end mit \lstinline[style=fileinline]|k_sem_init(&my_sem, initial_count, limit)| initialisiert. Es kann aber auch mit \lstinline[style=fileinline]|K_SEM_DEFINE(name, initial_count, count_limit)| zur {\"U}bersetzungszeit definiert und initialisiert werden. Das Geben ist durch \lstinline[style=fileinline]|k_sem_give(&my_sem)| m{\"o}glich. Das Nehmen funktioniert auf der anderen Seite durch \lstinline[style=fileinline]|k_sem_take(&my_sem, timeout)|, wobei \lstinline[style=fileinline]|timeout| als Wartezeit, z.B. \lstinline[style=fileinline]|K_MSEC(50)|, oder als einer der Sonderwerte \lstinline[style=fileinline]|K_NO_WAIT|/ \lstinline[style=fileinline]|K_FOREVER| angegeben werden kann.

\textbf{Empfohlene Einsatzgebiete }Die Zephyr-Dokumentation empfiehlt Semaphoren zur Zugriffskontrolle auf eine Menge von Ressourcen durch mehrere Threads und zur Synchronisation der Verarbeitung zwischen erzeugenden und verbrauchenden Threads oder ISRs.

\subsubsection{Interrupts}
Eine Interrupt-Service-Routine (ISR) {\"u}bernimmt die Bearbeitung von einer Unterbrechungsanforderung au{\ss}erhalb des regul{\"a}ren Thread-Ablaufs. Daf{\"u}r wird eine Interruptquelle mit einem Handler verbunden. Bei der Registrierung werden unter anderem die IRQ-Nummer, die Priorit{\"a}t, die auszuf{\"u}hrende Funktion und ein Argument festgelegt, das ggf. {\"u}bergeben werden soll (Zephyr Project, 2024, \glqq{}Interrupts\grqq{}, Abschnitt \glqq{}Concepts\grqq{}).

Die Zuordnung wird {\"u}ber die Interrupt-Infrastruktur des Systems gemacht. Interrupts, die nicht eingerichtet sind, werden dabei durch einen Standardhandler behandelt, der einen unerwarteten Aufruf als Fehler meldet. Zephyr unterst{\"u}tzt au{\ss}erdem gemeinsam genutzte Interruptleitungen. F{\"u}r diese Untersuchung wird jedoch eine einzelne, dedizierte Zuordnung verwendet (Zephyr Project, 2024, \glqq{}Interrupts\grqq{}, Abschnitte \glqq{}Concepts\grqq{} und \glqq{}Sharing interrupt lines\grqq{}).

F{\"u}r die Messung in Kapitel 5.2 ist insbesondere der Weg von der ausgel{\"o}sten IRQ-Anforderung bis zum Zeitstempel innerhalb des registrierten Handlers relevant.

\textbf{Interrupt-Verschachtelung }Der Kernel unterst{\"u}tzt das Verschachteln von Interrupts. Eine ISR kann dabei mitten in ihrer Ausf{\"u}hrung von einer h{\"o}herprioren ISR unterbrochen werden. Dabei setzt die niedrigerpriore ISR ihre Verarbeitung fort, sobald die h{\"o}herpriore fertig ist. Eine ISR wird dabei im sogenannten \glqq{}Interrupt-Kontext\grqq{} des Kernels ausgef{\"u}hrt, der {\"u}ber einen eigenen dedizierten Stack-Bereich verf{\"u}gt. Laut der Zephyr-Dokumentation d{\"u}rfen viele Kernel-APIs nur von Threads und nicht von ISRs aufgerufen werden. Jedoch stellt der Kernel f{\"u}r Routinen, die von beiden aufgerufen werden k{\"o}nnen, die Funktion \lstinline[style=fileinline]|k_is_in_isr()| bereit. Damit kann eine Routine ihr Verhalten basierend auf dem Aufrufkontext anpassen.

\textbf{Unterbrechungen verhindern }Ein Thread kann Interrupt-Verarbeitungen im System tempor{\"a}r {\"u}ber eine IRQ-Sperre (\lstinline[style=fileinline]|irq_lock()|/ \lstinline[style=fileinline]|irq_unlock()|) unterbinden. Diese Sperre ist laut der Zephyr-Dokumentation thread-spezifisch. Auf der anderen Seite kann ein einzelnes Interrupt gezielt {\"u}ber \lstinline[style=fileinline]|irq_disable()|/ \lstinline[style=fileinline]|irq_enable()| deaktiviert bzw. aktiviert werden. Dadurch wird die zugeh{\"o}rige ISR w{\"a}hrenddessen nicht ausgef{\"u}hrt.

\textbf{Auslagern von ISR-Arbeit }Um ein vorhersagbares (deterministisches) System sicherstellen zu k{\"o}nnen, muss eine ISR schnell ausgef{\"u}hrt werden. Daher empfiehlt die Zephyr-Dokumentation, zeitaufwendige Verarbeitungen an Threads auszulagern. Dies ist bspw. m{\"o}glich, indem die ISR {\"u}ber ein Kernel-Objekt wie ein Semaphor einen Hilfs-Thread benachrichtigt.

\textbf{Definition einer ISR Definition einer ISR:} Statische Interrupt-Verbindungen werden mit \lstinline[style=fileinline]|IRQ_CONNECT(irq, priority, isr, arg, flags)| zur {\"U}bersetzungszeit festgelegt. IRQ-Nummer, Priorit{\"a}t und Flags m{\"u}ssen dabei konstante Ausdr{\"u}cke sein. Anschlie{\ss}end wird die Leitung mit \lstinline[style=fileinline]|irq_enable()| aktiviert. F{\"u}r dynamische Verbindungen stellt Zephyr getrennte APIs bereit (Zephyr Project, 2024, \glqq{}Interrupts\grqq{}).

\textbf{Empfohlene Einsatzgebiete und Konfiguration }Die Zephyr-Dokumentation empfiehlt regul{\"a}re bzw. \glqq{}direkte\grqq{} ISRs f{\"u}r Interrupt-Verarbeitung, die eine schnelle, nicht-blockierende Reaktion erfordert. Dabei sollen zeitaufwendige oder blockierende Verarbeitungen an einen Thread {\"u}bergeben werden. Zudem kann man die Konfigurationsoption \lstinline[style=fileinline]|CONFIG_ISR_STACK_SIZE| hinzuf{\"u}gen, welche die Gr{\"o}{\ss}e des Interrupt-Stacks festlegt.

\subsubsection{Timer und Tickless-Kernel}
\textbf{Der Timer als Kernel-Objekt }Ein Timer ist laut Zephyr-Dokumentation ein Kernel-Objekt, welches den Zeitablauf anhand der Systemuhr des Kernels misst. Es k{\"o}nnen dabei mehrere Timer angelegt werden, wobei die Menge nur durch den verf{\"u}gbaren RAM begrenzt wird. Ein Zephyr-Timer kann einen einmaligen oder wiederkehrenden Ablauf ausl{\"o}sen. Beim Start werden daf{\"u}r zwei Zeitangaben festgelegt, darunter die Dauer, welche den ersten Ablauf bestimmt und die Periode, welche das Zeitraster der folgenden Abl{\"a}ufe bestimmt. Mit \lstinline[style=fileinline]|K_NO_WAIT| oder \lstinline[style=fileinline]|K_FOREVER| als Periodenparameter wird der Timer dabei einmalig betrieben. Eine zugeordnete Ablauffunktion wird dann im Kontext des Systemclock-Interrupt-Handlers aufgerufen (Zephyr Project, 2024, \glqq{}Timers\grqq{}, Abschnitt \glqq{}Concepts\grqq{}).

F{\"u}r den Benchmark wird ein periodischer Timer benutzt. Dabei ist die Unterscheidung zwischen seinem vorgesehenen Ablaufzeitpunkt und dem Eintritt in die Ablauffunktion f{\"u}r die Interpretation der Zeitmessung relevant.

\textbf{Wichtige Pr{\"a}zisierung f{\"u}r Zeitmessungen }Die Zephyr-Dokumentation erw{\"a}hnt, dass Dauer und Periode eines Timers Mindest-Verz{\"o}gerungen festlegen. Das liegt daran, dass durch die interne Pr{\"a}zision des Systemzeitgebers und andere m{\"o}gliche Laufzeiteinfl{\"u}sse wie Interrupt-Verz{\"o}gerungen mehr Zeit vergehen kann, als von den zugeh{\"o}rigen System-Zeit-APIs gemessen wird. Die angegebene Dauer ist eine Mindestverz{\"o}gerung bis zum Ablauf. Daraus folgt jedoch nicht, dass jedes zwischen zwei Callback-Eintritten gemessene Intervall nur nach oben abweichen kann, weil Verz{\"o}gerungen einzelner Callback-Aufrufe auch die benachbarten Intervalle beeinflussen k{\"o}nnen (Zephyr Project, 2024, \glqq{}Timers\grqq{}).

\textbf{Definition und Start eines Timers.} Ein Timer wird {\"u}ber eine Variable vom Typ \lstinline[style=fileinline]|k_timer| angelegt und mit \lstinline[style=fileinline]|k_timer_init()| initialisiert. Auf der anderen Seite kann dies schon zur {\"U}bersetzungszeit durch \lstinline[style=fileinline]|K_TIMER_DEFINE(name, expiry_fn, stop_fn)| definiert werden. Ein Timer wird {\"u}ber \lstinline[style=fileinline]|k_timer_start(timer, duration, period)| gestartet. Dabei wird sein Statuswert auf Null zur{\"u}ckgesetzt und der Timer f{\"a}ngt an, herunterzuz{\"a}hlen. Die Ablauffunktion (\emph{expiry function}), welche dem Timer zugeordnet wird, wird laut der Zephyr-Dokumentation vom Interrupt-Handler der Systemuhr selbst ausgef{\"u}hrt. Dies passiert jedes Mal, wenn der Timer abl{\"a}uft.

\textbf{Zeiteinheiten im Kernel.} Die Zephyr-Dokumentation erw{\"a}hnt verschiedene interne Zeiteinheiten:

\begin{itemize}
\item Reale Zeitwerte (Millisekunden/ Mikrosekunden) als portable Standarddarstellung
\item Einen \glqq{}Cycle"-Z{\"a}hler {\"u}ber \lstinline[style=fileinline]|k_cycle_get_32()|/ \lstinline[style=fileinline]|k_cycle_get_64()|, der den schnellsten verf{\"u}gbaren Zyklenz{\"a}hler des Systems darstellt
\item Sogenannte \glqq{}Ticks" als die interne Z{\"a}hleinheit, in der der Kernel sein gesamtes Uptime- und Timeout-Bookkeeping durchf{\"u}hrt
\end{itemize}
Kernel-Timeouts und Uptime werden intern in Ticks verwaltet. Deren Frequenz wird {\"u}ber \lstinline[style=fileinline]|CONFIG_SYS_CLOCK_TICKS_PER_SEC| festgelegt (Zephyr Project, 2024, \glqq{}Kernel Timing\grqq{}). Die Tick-Rate ist dabei {\"u}ber \lstinline[style=fileinline]|CONFIG_SYS_CLOCK_TICKS_PER_|SEC konfigurierbar. Die Uptime des Systems l{\"a}sst sich zudem {\"u}ber \lstinline[style=fileinline]|k_uptime_get()| in Millisekunden abfragen. {\"U}ber \lstinline[style=fileinline]|k_uptime_ticks()| kann die interne 64-Bit-Tick-Z{\"a}hlung abgefragt werden. Innerhalb der Verarbeitung eines Timeout-Ereignisses entspricht dabei dieser Wert in Zephyr v3.7.2 jedoch dem geplanten Timeout-Zeitpunkt und nicht unbedingt dem Eintrittszeitpunkt in die Ablauffunktion. F{\"u}r die Messung der realen Callback-Abst{\"a}nde wird deshalb in dieser Arbeit der frei laufende Systemtimer {\"u}ber \lstinline[style=fileinline]|sys_clock_cycle_get_32()| ausgelesen.

\textbf{Der tickless Kernel.} Ein Zeitgebertreiber muss dem Kernel neue Ticks {\"u}ber \lstinline[style=fileinline]|sys_clock_announce()| melden. Die Zephyr-Dokumentation beschreibt dabei zwei Implementierungsarten. Eine \glqq{}nat{\"u}rliche" Umsetzung dieser Schnittstelle f{\"u}hrt zu einem \glqq{}tickless\grqq{} Kernel, der Timer-Interrupts nur f{\"u}r registrierte Ereignisse empf{\"a}ngt und verarbeitet. Hierbei wird dies auf programmierbare Hardware-Z{\"a}hler gest{\"u}tzt, die unregelm{\"a}{\ss}ige Interrupts ausl{\"o}sen k{\"o}nnen. Auf der anderen Seite steht eine traditionelle, \glqq{}getaktete" Umsetzung, bei der der Treiber Interrupts in einer regelm{\"a}{\ss}igen Frequenz empf{\"a}ngt, die der Systemtaktrate entspricht. Zudem wird jedes Mal genau ein Tick gemeldet, wobei dies passiert, wenn ein Ereignis ansteht oder auch nicht.

\textbf{Empfohlene Einsatzgebiete.} Ein Timer soll laut Dokumentation verwendet werden, um einen asynchronen Vorgang nach einer bestimmten Zeit auszul{\"o}sen, oder um festzustellen, ob eine bestimmte Zeit bereits verstrichen ist. Das gilt hierbei besonders dann, wenn eine h{\"o}here Pr{\"a}zision ben{\"o}tigt wird, als es die Aufrufe \lstinline[style=fileinline]|k_sleep()| und \lstinline[style=fileinline]|k_usleep()| bieten. Auf der anderen Seite empfiehlt die Zephyr-Dokumentation, f{\"u}r die Messung der f{\"u}r einen Vorgang selbst ben{\"o}tigten Zeit direkt die System- bzw. Hardware-Uhr auszulesen, anstatt daf{\"u}r einen Timer zu benutzen.

\subsubsection{Speicherverwaltung}
Laut Zephyr-Dokumentation werden von Zephyr mehrere Werkzeuge bereitgestellt, mit denen Threads dynamisch Speicher anfordern k{\"o}nnen. F{\"u}r diese Arbeit sind davon zwei Ebenen relevant:

\begin{itemize}
\item Der synchronisierte \lstinline[style=fileinline]|k_heap|, den die Anwendung selbst f{\"u}r einen eigenen Speicherbereich anlegen kann
\item Der vordefinierte System-Heap, der {\"u}ber \lstinline[style=fileinline]|k_malloc()|/ \lstinline[style=fileinline]|k_free()| angesprochen wird
\end{itemize}
\textbf{Der synchronisierte Heap (}\lstinline[style=fileinline]|k_heap|\textbf{).} Der einfachste Weg, einen Heap anzulegen, ist laut Zephyr-Dokumentation die statische Definition {\"u}ber das Makro \lstinline[style=fileinline]|K_HEAP_DEFINE(name, bytes)|. Dieses erzeugt eine statische \lstinline[style=fileinline]|k_heap|-Variable, die einen Speicherbereich der angegebenen Gr{\"o}{\ss}e verwaltet. Speicher wird dabei {\"u}ber \lstinline[style=fileinline]|k_heap_alloc(heap, bytes, timeout)| angefordert. Die Funktion verh{\"a}lt sich hierbei {\"a}hnlich wie das \lstinline[style=fileinline]|malloc()| aus der Standard-C-Bibliothek und liefert bei einem Fehler einen Nullzeiger (nullpointer) zur{\"u}ck. Der Heap unterst{\"u}tzt dabei auch blockierende Operationen. Ein Thread kann in den Schlafzustand versetzt werden, bis Speicher verf{\"u}gbar ist. Dabei gibt der Timeout-Parameter an, wie lange maximal gewartet werden darf. Freigegeben wird zuvor angeforderter Speicher {\"u}ber \lstinline[style=fileinline]|k_heap_free()|, genau wie bei \lstinline[style=fileinline]|free()|.

\textbf{Die Implementierung (}\lstinline[style=fileinline]|sys_heap|\textbf{)} Die Speicherzuteilung eines \lstinline[style=fileinline]|k_heap| basiert auf dem \lstinline[style=fileinline]|sys_heap-Allokator|. Dabei erg{\"a}nzt \lstinline[style=fileinline]|k_heap| die ben{\"o}tigte Synchronisation. Bei direkter Nutzung von \lstinline[style=fileinline]|sys_heap| muss der Aufrufer miteinander konkurrierende Zugriffe selbst organisieren (Zephyr Project, 2024, \glqq{}Memory Heaps\grqq{}, Abschnitt \glqq{}Low Level Heap Allocator\grqq{}).

Der Allokator verwaltet den Speicher in Einheiten von acht Byte und h{\"a}lt Informationen {\"u}ber belegte und freie Bereiche in Verwaltungsstrukturen. Freie Bl{\"o}cke werden dabei nach Gr{\"o}{\ss}enklassen organisiert, damit f{\"u}r eine Anforderung ein passender Block gesucht werden kann. Bei einer Freigabe k{\"o}nnen benachbarte freie Bereiche zusammengef{\"u}hrt werden. Diese Verfahren sorgend daf{\"u}r, dass Fragmentierung verringert wird. Jedoch schlie{\ss}en sie dies nicht aus (Zephyr Project, 2024, \glqq{}Memory Heaps\grqq{}, Abschnitt \glqq{}Low Level Heap Allocator\grqq{}).

Die Suche innerhalb einer Gr{\"o}{\ss}enklasse ist durch \lstinline[style=fileinline]|CONFIG_SYS_HEAP_ALLOC_LOOPS| begrenzt. Der dokumentierte Standardwert betr{\"a}gt hierbei drei. Dadurch wird der Suchaufwand begrenzt, jedoch kann die Auswahl eines passenderen Blocks darunter leiden. Ein begrenzter algorithmischer Aufwand wird dabei von einer konstanten, bei jedem Aufruf gemessenen Zyklenzahl unterschieden (Zephyr Project, 2024, \glqq{}Memory Heaps\grqq{}, Abschnitt \glqq{}Low Level Heap Allocator\grqq{}).

\textbf{Der System-Heap (}\lstinline[style=fileinline]|k_malloc|\textbf{/ }\lstinline[style=fileinline]|k_free|\textbf{).} Der System-Heap ist ein vordefinierter Speicherzuteiler, aus dem Threads in einer \lstinline[style=fileinline]|malloc()|-{\"a}hnlichen Weise dynamisch Speicher aus einem gemeinsamen Speicherbereich anfordern k{\"o}nnen. Anders als bei anderen Heaps kann der System-Heap nicht direkt {\"u}ber eine Speicheradresse referenziert werden. Es gibt laut Zephyr-Dokumentation nur einen einzigen System-Heap im gesamten System. Seine Gr{\"o}{\ss}e wird {\"u}ber die Kconfig-Option \lstinline[style=fileinline]|CONFIG_HEAP_MEM_POOL_SIZE| festgelegt. Dabei ist der Standardwert Null Byte, was bedeutet, dass der Kernel gar kein Speicherpool-Objekt anlegt. Wenn die angeforderte Gr{\"o}{\ss}e hierbei gr{\"o}{\ss}er ist als der verf{\"u}gbare Speicher, dann schl{\"a}gt der Build schon in der Linker-Phase fehl. Speicher wird {\"u}ber \lstinline[style=fileinline]|k_malloc(size)| angefordert und {\"u}ber \lstinline[style=fileinline]|k_free(ptr)| wieder freigegeben. Zudem entsprechen beide Funktionen, was die Semantik angeht, den gleichnamigen Standard-C-Funktionen.

\subsubsection{Logging-Subsystem}
Die Logging-API von Zephyr stellt laut Zephyr-Dokumentation eine Schnittstelle bereit, {\"u}ber die von Entwicklern ausgegebene Nachrichten verarbeitet werden. Nachrichten laufen dabei zuerst durch ein Frontend und werden dann von aktiven Backends verarbeitet. Die Architektur des Logging-Subsystems besteht dabei laut Zephyr-Dokumentation aus genau drei Hauptteilen, dem Frontend, dem Core und den Backends.

\textbf{Das Standard-Frontend} Das Standard-Frontend wird immer dann aktiv, wenn die Logging-API in einer Quelle aufgerufen wird, z.B. durch das Makro \lstinline[style=fileinline]|LOG_INF|. Es ist daf{\"u}r zust{\"a}ndig, eine Nachricht zu filtern, sowohl zur {\"U}bersetzungs- als auch zur Laufzeit. Au{\ss}erdem ist es daf{\"u}r zust{\"a}ndig, Speicher f{\"u}r die Nachricht zu reservieren, die Nachricht zu erzeugen und sie dann zu {\"u}bergeben. Da die Logging-API laut Zephyr-Dokumentation auch innerhalb eines Interrupts aufgerufen werden kann, ist das Frontend darauf optimiert, eine Nachricht so schnell wie m{\"o}glich zu protokollieren. Eine Log-Nachricht wird dabei in einem zusammenh{\"a}ngenden Speicherblock abgelegt, der aus einem \glqq{}zirkul{\"a}ren Paketpuffer\grqq{} reserviert wird.

\textbf{Zwei }Die Zephyr-Dokumentation unterscheidet drei globale Logging-Modi, darunter Deferred, Immediate und den auf minimalen Funktionsumfang ausgerichteten Minimal-Modus (Zephyr Project, 2024, \glqq{}Logging\grqq{}). Die Zusammenfassung der Logging-Funktionen auf derselben Zephyr-Dokumentationsseite notiert dazu, dass \glqq{}deferred logging" die f{\"u}r eine Log-Ausgabe ben{\"o}tigte Zeit reduziert. Dies wird m{\"o}glich, indem zeitintensive Operationen in einen bekannten, daf{\"u}r vorgesehenen Kontext verschoben werden, anstatt die Nachricht schon beim eigentlichen Aufruf zu verarbeiten und auszugeben.

\textbf{Auswirkung auf den Stack-Verbrauch} Die Zephyr-Dokumentation erl{\"a}utert im Abschnitt \glqq{}Stack usage\grqq{}, dass die Aktivierung des Logging-Subsystems den Stack-Verbrauch des jeweils aufrufenden Kontexts beeinflusst:

\begin{itemize}
\item Wird \lstinline[style=fileinline]|CONFIG_LOG_MODE_DEFERRED| verwendet, wird dieser Verbrauch geringer, da die Logging-Funktion an der Aufrufstelle nur die Nachricht erzeugt und im Puffer ablegt.
\item Wird \lstinline[style=fileinline]|CONFIG_LOG_MODE_IMMEDIATE| verwendet, wird die Nachricht direkt vom Backend verarbeitet, was auch die Zeichenketten-Formatierung inkludiert. Der Stack-Verbrauch h{\"a}ngt in diesem Fall auch davon ab, welche Backends verwendet werden.
\end{itemize}
\textbf{Verarbeitung im Deferred-Modus} Normalerweise wird die Verarbeitung von Nachrichten im Deferred-Modus intern durch einen daf{\"u}r vorgesehenen Task {\"u}bernommen, der automatisch startet. Dieser ist steuerbar {\"u}ber \lstinline[style=fileinline]|CONFIG_LOG_PROCESS_THREAD|. Die Aktivierungsschwelle l{\"a}sst sich {\"u}ber die Anzahl gepufferter Nachrichten durch \lstinline[style=fileinline]|CONFIG_LOG_PROCESS_TRIGGER_THRESHOLD| konfigurieren. Die Gr{\"o}{\ss}e des daf{\"u}r verwendeten Ringpuffers wird dabei {\"u}ber \lstinline[style=fileinline]|CONFIG_LOG_BUFFER_SIZE| festgelegt. {\"U}ber \lstinline[style=fileinline]|CONFIG_LOG_DEFAULT_LEVEL| l{\"a}sst sich auch der Schweregrad festlegen, der f{\"u}r Module ohne eigene Log-Level-Angabe benutzt wird. Die Zephyr-Dokumentation nennt vier verf{\"u}gbare Schweregrade. Diese sind Fehler (Error), Warnung (Warning), Information (Info) und Debug. Als Ausgabekanal (\emph{Backend}) steht unter anderem \lstinline[style=fileinline]|CONFIG_LOG_BACKEND_UART| zur Verf{\"u}gung. Laut Zephyr-Dokumentation unterst{\"u}tzt das Subsystem bis zu neun gleichzeitig aktive Backends. {\"U}ber \lstinline[style=fileinline]|CONFIG_LOG_FUNC_NAME_PREFIX_DBG| l{\"a}sst sich zudem festlegen, ob Debug-Nachrichten automatisch mit dem Namen der aufrufenden Funktion versehen werden.

\textbf{Vom Projekt selbst gemessene Gr{\"o}{\ss}enordnung} Die Dokumentation enth{\"a}lt im Abschnitt \glqq{}Benchmark" eigene, auf der Plattform qemu\_x86 gemessene Vergleichswerte des Zephyr-Projekts (Testsuite tests/subsys/logging/log\_benchmark), die den Zeitaufwand von einzelnen Logging-Operationen in der Gr{\"o}{\ss}enordnung von ca. 7 bis 50 \textmu{}s notieren, abh{\"a}ngig von Modus und Art der Daten.

\subsubsection{Timing-API}
Die Zephyr-Timing-API bietet Funktionen, mit denen die Ausf{\"u}hrungszeit eines Codeabschnitts herausgefunden werden kann, um diesen zu analysieren oder zu optimieren. Die Zephyr-Dokumentation erw{\"a}hnt dabei, dass die Timing-Funktionen einen anderen Zeitgeber benutzen k{\"o}nnen als der Standard-Kernel-Timer. Welcher Zeitgeber hierbei verwendet wird, ist architektur-, SoC- oder board-spezifisch festgelegt. Die API ist daher bewusst allgemein gehalten und empfiehlt keinen bestimmten Hardware-Z{\"a}hler. Die in dieser Arbeit verwendete Hardware-Plattform wird in Kapitel 4.2 genauer erkl{\"a}rt.

\textbf{Voraussetzung} Um die Timing-Funktionen nutzen zu k{\"o}nnen, muss die Kconfig-Option \lstinline[style=fileinline]|CONFIG_TIMING_FUNCTIONS| aktiviert sein.

\textbf{Vorgesehener Ablauf laut Dokumentation} Die Nutzung der API folgt einem festen Schema. Die Zeitmessung kann in Vorbereitung, Erfassung und Auswertung unterteilt werden. Zur Vorbereitung initialisieren \lstinline[style=fileinline]|timing_init()| und \lstinline[style=fileinline]|timing_start()| die benutzte Timing-Infrastruktur und starten ihre Nutzung. F{\"u}r einen Codeabschnitt werden dann anschlie{\ss}end mit \lstinline[style=fileinline]|timing_counter_get()| zwei Z{\"a}hlerst{\"a}nde erfasst. \lstinline[style=fileinline]|timing_cycles_get()| bestimmt daraus dann die vergangene Zyklenzahl. Wenn dies gebraucht wird, kann die Zyklenzahl mit \lstinline[style=fileinline]|timing_cycles_to_ns()| in Nanosekunden umgerechnet werden. Nach Abschluss der Erfassung kann \lstinline[style=fileinline]|timing_stop()| benutzt werden (Zephyr Project, 2023, \glqq{}Executing Time Functions\grqq{}, Abschnitt \glqq{}Usage\grqq{}).

Die Messbibliothek dieser Arbeit {\"u}bernimmt die Initialisierung gemeinsam f{\"u}r die Benchmarks und unterteilt die Z{\"a}hlerzugriffe in \lstinline[style=fileinline]|bench_now()| und \lstinline[style=fileinline]|bench_elapsed()|.

\textbf{Weitere Funktionen der API} Au{\ss}er den genannten Kernfunktionen stellt die API \lstinline[style=fileinline]|timing_freq_get()|, welches die Frequenz des verwendeten Z{\"a}hlers in Hz liefert und \lstinline[style=fileinline]|timing_freq_get_mhz()|, was dasselbe in MHz ist, zur Verf{\"u}gung. Die Zephyr-Dokumentation notiert au{\ss}erdem, dass bei manchen Architekturen oder SoCs die vom Z{\"a}hler gelieferten Rohwerte noch skaliert werden m{\"u}ssen, um die Zyklenzahl zu erhalten. Diese Umrechnung {\"u}bernimmt \lstinline[style=fileinline]|timing_cycles_get()| intern.

\subsection{Build-Zeit-Konzepte}
Zephyr stellt ein eigenes Build- und Konfigurationssystem bereit, welches festlegt, wie eine Anwendung aus Quellcode zu einer lauff{\"a}higen Firmware zusammengesetzt wird. Die Zephyr-Dokumentation fasst darunter das CMake-basierte Build-System, den Devicetree, das Kconfig-Konfigurationssystem, Snippets und das Sysbuild-System f{\"u}r Mehr-Image-Builds zusammen.

Das Konfigurationssystem Kconfig macht es m{\"o}glich, den Zephyr-Kernel und seine Subsysteme zur Build-Zeit an die Anforderungen einer bestimmten Anwendung und Zielplattform anzupassen. Dabei handelt es sich um dasselbe Konfigurationssystem, das auch der Linux-Kernel verwendet. Das Ziel ist dabei eine Konfigurierbarkeit, ohne dass daf{\"u}r Quellcode ver{\"a}ndert werden muss. Hierbei werden Konfigurationsoptionen (auch \glqq{}Symbole\grqq{} genannt) in \lstinline[style=fileinline]|Kconfig|-Dateien definiert. Diese legen dabei gleichzeitig die Abh{\"a}ngigkeiten zwischen den Symbolen fest, die eine g{\"u}ltige Konfiguration bestimmen. Aus Kconfig entsteht als Ergebnis eine Header-Datei mit Makros, die zur Build-Zeit ausgewertet werden k{\"o}nnen. Dadurch kann nicht ben{\"o}tigter Code aus dem Build ausgeschlossen werden.

Damit auch externe Software-Projekte in dieses Build-System eingebunden werden k{\"o}nnen, ohne deren Code direkt in den Zephyr-Quellbaum zu kopieren, verwendet Zephyr das Konzept der Module. Dabei handelt es sich um Repositories, die eine eigene \lstinline[style=fileinline]|module.yml|-Datei mitbringen und dadurch vom Build-System automatisch mit ihren eigenen \lstinline[style=fileinline]|CMakeLists.txt|- und \lstinline[style=fileinline]|Kconfig|-Dateien eingebunden werden.

Das Kommandenzeilenwerkzeug West {\"u}bernimmt dabei die Verwaltung mehrerer, teilweise unabh{\"a}ngiger Repositories, darunter Zephyr selbst, Module und eigene Anwendungsprojekte. Die Zephyr-Dokumentation bezeichnet dieses Werkzeug auch als eine Art \glqq{}Schweizer Taschenmesser\grqq{} f{\"u}r das Projekt. West stellt hierbei ein Mehr-Repository-Verwaltungssystem bereit, das von Googles Repo-Werkzeug und Git-Submodulen inspiriert ist. Dabei ist West zudem erweiterbar. Zephyr selbst nutzt diese Erweiterbarkeit, um weitere Kommandos f{\"u}r das Bauen, Flashen und Debuggen von Anwendungen bereitzustellen.

F{\"u}r diese Arbeit sind dabei vor allem die Funktionsweise von Kconfig-Fragmenten, der Modul-Mechanismus zur Einbindung von eigenen Bibliotheken, West als Werkzeug zur Versionsverwaltung und zum Bauen/ Flashen und zuletzt die im Build-System enthaltenen Werkzeuge zur Analyse des Speicherbedarfs relevant.

\subsubsection{Kconfig}
\textbf{Grundprinzip} Kconfig legt fest, welche konfigurierbaren Softwarefunktionen in einen Zephyr-Build eingehen. Dazu definiert das Projekt Symbole mit ihren m{\"o}glichen Werten und Abh{\"a}ngigkeiten. Die Anwendung weist dann diesen Symbolen Werte zu, ohne dass daf{\"u}r die Implementierung der jeweiligen Systemfunktion ver{\"a}ndert werden muss. Aus der aufgel{\"o}sten Konfiguration entstehen dann anschlie{\ss}end unter anderem Makros, auf die der {\"U}bersetzungsprozess achten kann (Zephyr Project, 2024, \glqq{}Setting Kconfig configuration values\grqq{}).

F{\"u}r den Versuchsaufbau ist vor allem die Trennung zwischen gemeinsamer Anwendungskonfiguration und weiteren Fragmenten relevant. \lstinline[style=fileinline]|CONF_FILE| legt dabei die benutzte Anwendungskonfiguration fest. \lstinline[style=fileinline]|EXTRA_CONF_FILE| erweitert weitere Konfigurationsdateien. Auf diesem Mechanismus basiert die Variation der Profile in Kapitel 4.4 (Zephyr Project, o. D., \glqq{}Application Development\grqq{}, Abschnitt \glqq{}Important Build System Variables\grqq{}).

\textbf{Sichtbare und unsichtbare Symbole} Die Zephyr-Dokumentation unterteilt die Konfigurationsoptionen in \glqq{}sichtbare\grqq{} und \glqq{}unsichtbare\grqq{} Symbole. Sichtbare Symbole sind {\"u}ber einen Prompt-Text definiert, erscheinen in interaktiven Konfigurationswerkzeugen und k{\"o}nnen in Konfigurationsdateien gesetzt werden. Unsichtbare Symbole werden ohne einen Prompt definiert. Deren Wert ergibt sich nur aus Vorgabewerten oder {\"u}ber \lstinline[style=fileinline]|select|-Beziehungen anderer Symbole.

\textbf{Zuweisung in Konfigurationsdateien} Sichtbare Symbole werden in Konfigurationsdateien {\"u}ber \lstinline[style=fileinline]|CONFIG_<Symbolname>=<Wert>| gesetzt. Dabei d{\"u}rfen um das Gleichheitszeichen keine Leerzeichen stehen. Boolesche Symbole werden zudem {\"u}ber \lstinline[style=fileinline]|y| bzw. \lstinline[style=fileinline]|n| aktiviert oder deaktiviert. Ein auf \lstinline[style=fileinline]|n| gesetztes Symbol kann auf der anderen Seite auch als Kommentar in der Form \lstinline[style=fileinline]|# CONFIG_<Symbolname> is not set| geschrieben werden. Genau dieses Format befindet sich auch in der zusammengef{\"u}hrten Konfigurationsdatei \lstinline[style=fileinline]|zephyr/.config|.

\textbf{Die initiale Konfiguration} Laut der Zephyr-Dokumentation entsteht die initiale Konfiguration einer Anwendung durch Zusammenf{\"u}hren von drei Quellen:

\begin{itemize}
\item Die board-spezifische \lstinline[style=fileinline]|<BOARD>_defconfig|-Datei
\item Die CMake-Cache-Eintr{\"a}ge, die mit \lstinline[style=fileinline]|CONFIG_| beginnen
\item Die Anwendungskonfiguration selbst, standardm{\"a}{\ss}ig \lstinline[style=fileinline]|prj.conf|
\end{itemize}
Ist anstelle von \lstinline[style=fileinline]|prj.conf| die Variable \lstinline[style=fileinline]|CONF_FILE| gesetzt, werden die dort angegebenen Konfigurationsdateien zusammengef{\"u}hrt und als Anwendungskonfiguration benutzt. Jede Datei dieser Art wird dabei als \glqq{}Kconfig\emph{-}Fragment\grqq{} bezeichnet. \lstinline[style=fileinline]|CONF_FILE| kann unter anderem direkt {\"u}ber \lstinline[style=fileinline]|-DCONF_FILE=<Datei(en)>| beim Aufruf von \lstinline[style=fileinline]|west build| gesetzt werden. Wird ein Symbol sowohl in der \lstinline[style=fileinline]|<BOARD>_defconfig| als auch in der Anwendungskonfiguration zugewiesen, hat laut Zephyr-Dokumentation der  Wert Vorrang, der in der Anwendungskonfiguration gesetzt wurde. Das Ergebnis der Zusammenf{\"u}hrung wird als \lstinline[style=fileinline]|zephyr/.config| im Build-Verzeichnis gespeichert.

\textbf{Konfigurieren eines choice} Ein choice, also mehrere Symbole, von denen laut Zephyr-Dokumentation zu jedem Zeitpunkt genau eines aktiv sein kann, l{\"a}sst sich konfigurieren, indem eines seiner Symbole in einer Konfigurationsdatei auf \lstinline[style=fileinline]|y| gesetzt wird. Dadurch werden automatisch alle restlichen Symbole des choice auf \lstinline[style=fileinline]|n| gesetzt. Wenn mehrere Symbole vom selben choice auf \lstinline[style=fileinline]|y| gesetzt werden, dann wird laut Zephyr-Dokumentation nur das zuletzt gesetzte ber{\"u}cksichtigt, w{\"a}hrend die restlichen auf \lstinline[style=fileinline]|n| zur{\"u}ckgesetzt werden.

\subsubsection{CMake und Zephyr-Module}
\textbf{Build-System} CMake beschreibt den Aufbau der Firmware {\"u}ber Targets und deren Abh{\"a}ngigkeiten. F{\"u}r Anwendungscode stellt Zephyr dabei vor allem das Bibliotheks-Target \lstinline[style=fileinline]|app| bereit. Die dazugeh{\"o}rigen Quelldateien werden bspw. mit \lstinline[style=fileinline]|target_sources(app PRIVATE ...)| zugeordnet. \lstinline[style=fileinline]|PRIVATE| grenzt dabei diese Quellenzuordnung auf das betreffende Target ein. Mit \lstinline[style=fileinline]|PUBLIC| kann sie zudem Teil seiner Schnittstelle f{\"u}r abh{\"a}ngige Targets werden (Zephyr Project, 2024, \glqq{}Build System (CMake)\grqq{}).

Aus diesen Beschreibungen erzeugt CMake zuerst die Bauvorschriften f{\"u}r das ausgew{\"a}hlte Host-Werkzeug, bspw. Ninja oder Make. Dieses f{\"u}hrt dann anschlie{\ss}end die {\"U}bersetzungs- und Linkschritte aus. Dadurch sind die Konfiguration des Bauprozesses und seine eigentliche Ausf{\"u}hrung voneinander getrennt (Zephyr Project, 2024, \glqq{}Build System (CMake)\grqq{}, Abschnitt \glqq{}Configuration Phase\grqq{}).

\textbf{Externe Projekte als Module} Um auf den Quellcode externer Projekte zur{\"u}ckzugreifen, benutzt Zephyr das Konzept der \glqq{}Module\grqq{}. Ein Modul muss daf{\"u}r {\"u}ber eine Datei \lstinline[style=fileinline]|module.yml| verf{\"u}gen, durch welche das Zephyr-Build-System den Quellcode aus dem betreffenden Repository einbindet. Laut Zephyr-Dokumentation soll ein Modul dabei kein Code sein, der nur f{\"u}r Zephyr geschrieben wurde. Solcher Code sollte direkt in den Zephyr-Baum eingebaut werden.

\textbf{Technische Einbindung} Die Build-System-Variable \lstinline[style=fileinline]|ZEPHYR_MODULES| ist eine CMake-Liste absoluter Pfade zu den Verzeichnissen, die Zephyr-Module enthalten. Diese Module enthalten \lstinline[style=fileinline]|CMakeLists.txt|- und \lstinline[style=fileinline]|Kconfig|-Dateien, welche beschreiben, wie sie gebaut bzw. konfiguriert werden. Die \lstinline[style=fileinline]|CMakeLists.txt|-Dateien der Module werden dabei dem Build {\"u}ber CMakes \lstinline[style=fileinline]|add_subdirectory()|-Befehl hinzugef{\"u}gt. Die \lstinline[style=fileinline]|Kconfig|-Dateien werden in den Kconfig-Men{\"u}baum eingebunden. Wenn West installiert und \lstinline[style=fileinline]|ZEPHYR_MODULES| nicht gesetzt ist, dann findet das Build-System die Module selbst. Es ruft dabei \lstinline[style=fileinline]|west list| auf, um die Pfade aller Projekte im West-Arbeitsbereich zu erhalten. Anschlie{\ss}end pr{\"u}ft es f{\"u}r jedes Projekt, ob eine \lstinline[style=fileinline]|module.yml|-Datei vorhanden ist. Wenn dem so ist, dann wird ihr Inhalt zur Einbindung verwendet. Fehlt jedoch diese Datei, dann pr{\"u}ft das Build-System, ob im Projekt eine \lstinline[style=fileinline]|zephyr/CMakeLists.txt| und eine \lstinline[style=fileinline]|zephyr/Kconfig| vorhanden sind. Wenn beide vorhanden sind, dann gilt das Projekt ebenfalls als Modul.

\subsubsection{West}
West hat im Zephyr-Umfeld zwei miteinander verbundene Aufgaben. Seine eingebauten Kommandos verwalten die Repositories, die zu einem Arbeitsbereich geh{\"o}ren und deren im Manifest festgelegte Revisionen. Zudem erg{\"a}nzen Erweiterungskommandos projektspezifische Funktionen. Zephyr benutzt diesen Mechanismus unter anderem f{\"u}r das Bauen, Flashen und Debuggen von Anwendungen (Zephyr Project, 2023, \glqq{}West (Zephyr\grq{}s meta-tool)\grqq{}).

F{\"u}r diese Untersuchung ist vor allem wichtig, dass die Zephyr-Version und die dazugeh{\"o}rigen Module zusammen {\"u}ber das Projektmanifest festgelegt werden k{\"o}nnen. Dadurch l{\"a}sst sich der Softwarestand nachvollziehbar beschreiben, der f{\"u}r eine Messreihe vorgesehen ist.

\textbf{Aufrufsyntax} Wie bei git oder docker lautet ein Aufruf von West: \lstinline[style=fileinline]|west [allgemeine-Optionen] <Kommando> [Optionen] <Argumente>|. Seit Version 0.8 l{\"a}sst sich West auch {\"u}ber \lstinline[style=fileinline]|python3 -m west [...]| aufrufen. {\"U}ber \lstinline[style=fileinline]|west --help| bzw. \lstinline[style=fileinline]|west <Kommando> -h| zeigt das Werkzeug Hilfetexte zu den verf{\"u}gbaren bzw. zu einem bestimmten Kommando.

\textbf{Unterst{\"u}tzte Topologien} Die West-Dokumentation beschreibt drei unterst{\"u}tzte Arbeitsbereich-Topologien:

\begin{itemize}
\item \textbf{T1}: \glqq{}Star"-Topologie. Zephyr selbst ist das Manifest-Repository. Dies ist die Standardeinstellung, analog zu Git-Submodulen mit Zephyr als Super-Projekt
\item \textbf{T2}: \glqq{}Star"-Topologie. Eine Zephyr-Anwendung ist das Manifest-Repository, was n{\"u}tzlich ist f{\"u}r Projekte, die sich auf eine einzelne Anwendung konzentrieren
\item \textbf{T3}: Wald-Topologie mit einem eigenen, quellcodefreien Manifest-Repository f{\"u}r mehrere unabh{\"a}ngige Anwendungen
\end{itemize}
\textbf{T2-Topologie als Projektaufbau} Ein Repository, das eine Zephyr-Anwendung enth{\"a}lt, fungiert als zentrales Repository und nennt in seiner eigenen \lstinline[style=fileinline]|west.yml| die weiteren f{\"u}r den Build ben{\"o}tigten Projekte, inklusive des Zephyr-Repositorys selbst und aller Module. Die West-Dokumentation zeigt dazu ein Beispiel-Manifest, das {\"u}ber den Schl{\"u}ssel \lstinline[style=fileinline]|import: true| beim Projekt \lstinline[style=fileinline]|zephyr| und den Schl{\"u}ssel \lstinline[style=fileinline]|self: path: application| genau dieses Muster umsetzt. Der Vorteil des import-Mechanismus liegt laut West-Dokumentation darin, dass die Revisionen der importierten Module nicht getrennt nachverfolgt werden m{\"u}ssen, sondern automatisch aus der Revision der \lstinline[style=fileinline]|zephyr/west.yml| bestimmt werden.

\textbf{Der manifest-rev-Branch} West legt in jedem Projekt einen Git-Branch namens \lstinline[style=fileinline]|manifest-rev| an und verwaltet diesen dann selbst. Dieser zeigt zudem auf die Revision, die das Manifest beim letzten Aufruf von \lstinline[style=fileinline]|west update| f{\"u}r das jeweilige Projekt angegeben hat. Die West-Dokumentation warnt jedoch davor, diesen Branch manuell zu ver{\"a}ndern, da West ihn beim n{\"a}chsten Update zur{\"u}cksetzt.

\textbf{Das Manifest-Kommando }{\"U}ber \lstinline[style=fileinline]|west manifest --path| kann der Pfad zur obersten Manifest-Datei ausgeben werden. Genau dieser Befehl wird zur Kontrolle nach der Einrichtung des Arbeitsbereichs benutzt.

\subsubsection{Build-System-Werkzeuge zur Analyse des Speicherbedarfs}
Das Zephyr-Build-System bietet in der, bei dieser Untersuchung verwendeten Version v3.7.2 (Zephyr-Dokumentation: v3.7.0) genau drei Targets zur Analyse des Speicherbedarfs: \lstinline[style=fileinline]|ram_report|, \lstinline[style=fileinline]|rom_report| und \lstinline[style=fileinline]|puncover|, welches um \lstinline[style=fileinline]|pahole| zur Analyse von Datenstrukturen erg{\"a}nzt wird.

\textbf{Grundprinzip} Das Build-System bietet laut Zephyr-Dokumentation drei Targets, um RAM-, ROM- und Stack-Nutzung in erzeugten Firmware-Abbildern zu betrachten und zu analysieren. Die Werkzeuge arbeiten auf dem fertigen Abbild und liefern Informationen {\"u}ber die Gr{\"o}{\ss}e von Symbolen und Code sowohl im RAM als auch im ROM. Symbole, die sich nicht anhand von Metadaten einer Quelldatei zuordnen lassen, werden dabei in zwei Sammelkategorien ausgewiesen. \lstinline[style=fileinline]|Hidden| bezeichnet ein Symbol ohne zugeordnete Datei. \lstinline[style=fileinline]|No paths| sagt aus, dass der Funktionsname bekannt, aber kein eindeutiger Pfad im Verzeichnisbaum auffindbar ist.

\textbf{Target }\lstinline[style=fileinline]|ram_report| Dieses Target listet alle kompilierten Objekte und ihre RAM-Nutzung in Form einer Tabelle auf, mit Byte-Angabe und Prozentanteil pro Symbol. Die Daten werden dabei basierend auf den Speicherort des Objekts im Verzeichnisbaum und der Datei, die das Symbol enth{\"a}lt, gruppiert. Aufgerufen wird das Target laut Zephyr-Dokumentation nach einem vorherigen Build {\"u}ber \lstinline[style=fileinline]|west build -t ram_report|.

\textbf{Target }\lstinline[style=fileinline]|rom_report| Dieses Target funktioniert analog zu \lstinline[style=fileinline]|ram_report|, listet jedoch die ROM-Nutzung (Flash-Speicher) statt der RAM-Nutzung auf. Der Aufruf l{\"a}uft {\"u}ber \lstinline[style=fileinline]|west build -t rom_report|.

\textbf{Target }\lstinline[style=fileinline]|puncover| Dieses Target verwendet das gleichnamige Drittanbieter-Werkzeug, das nach dem Bauen einen lokalen Webserver startet, {\"u}ber den sich die Dateien des Projekts, inklusive ihrer ROM-, RAM- und Stack-Nutzung, im Browser interaktiv durchsuchen lassen. F{\"u}r eine Worst-Case-Stack-Nutzungsanalyse muss daf{\"u}r noch \lstinline[style=fileinline]|CONFIG_STACK_USAGE=y| gesetzt werden. Das Werkzeug muss zudem vor der Nutzung getrennt per \lstinline[style=fileinline]|pip3| installiert werden und wird von der Zephyr-Dokumentation als Drittanbieter-Werkzeug mit m{\"o}glicherweise wechselnder Funktionsf{\"a}higkeit notiert.

\textbf{Werkzeug }\lstinline[style=fileinline]|pahole| Bei \lstinline[style=fileinline]|pahole| handelt es sich um ein Werkzeug zur Analyse von Objektdateien, das die Gr{\"o}{\ss}e von Datenstrukturen und, durch Alignment auf die Wortbreite der CPU entstehende \glqq{}L{\"o}cher" in diesen Strukturen zeigt. Es muss getrennt installiert werden, z. B. {\"u}ber das Paket \lstinline[style=fileinline]|dwarves| unter Ubuntu. Jedoch ist es f{\"u}r diese Untersuchung nicht weiter relevant, da hier keine Datenstruktur-Layout-Optimierung im Fokus steht.

\section{Benchmark Metriken}
In den vorherigen Abschnitten wurden die Grundlagen zu Echtzeitsystemen, eingebetteten Systemen und das in dieser Arbeit verwendeten Zephyr RTOS definiert. Zudem wurden Netzwerk- und Sicherheitsprotokolle erl{\"a}utert. In diesem Abschnitt werden die f{\"u}nf in dieser Untersuchung gemessenen benchmark-Metriken definiert. Diese Definitionen sind bewusst basierend auf der allgemeinen Literatur zu Echtzeitsystemen definiert, damit sie auch unabh{\"a}ngig von einer bestimmten RTOS-Implementierung g{\"u}ltig sind. Die Implementierung dieser Metriken innerhalb von Zephyr wird in Kapitel 4.5 und 5 erl{\"a}utert.

\subsection{Context Switch}
\textbf{Kontext \& Kontextwechsel }Eine unterbrochene Aufgabe muss bei ihrer sp{\"a}teren Fortsetzung auf ihren vorherigen Zustand zur{\"u}ckgreifen k{\"o}nnen. Zu diesem Kontext geh{\"o}ren dabei vor allem Registerinhalte, der Programmz{\"a}hler und ggf. weitere aufgabenspezifische Informationen zum Zustand. Diese werden bspw. auf einem Stack gesichert (Laplante \& Ovaska, 2012, S. 88--89).

Beim Wechsel zu einer anderen Aufgabe wird der Kontext gesichert und der f{\"u}r die n{\"a}chste Aufgabe gebrauchte Zustand wiederhergestellt. Dieser Verwaltungsvorgang wird Kontextwechsel genannt. Seine Dauer ist Teil der Antwortzeit, ohne die eigentliche Anwendungsaufgabe voranzubringen. Deshalb sollte die Sicherung auf die Informationen begrenzt werden, die f{\"u}r eine korrekte Fortsetzung gebraucht werden (Laplante \& Ovaska, 2012, S. 88--89).

Abschnitt 2.5.1 erl{\"a}utert, wann Zephyr einen anderen Thread ausw{\"a}hlen kann. Der allgemeine Kontextwechsel beinhaltet das Sichern des f{\"u}r die Fortsetzung gebrauchten Zustands des bisherigen Threads und das Wiederherstellen des Zustands des n{\"a}chsten Threads (Laplante \& Ovaska, 2012, S. 88--89; Zephyr Project, 2024, \glqq{}Scheduling\grqq{}).

\subsection{Interrupt Latency}
\textbf{Interrupt-Latenz }Zwischen einer Interrupt-Anforderung und dem Beginn der zugeh{\"o}rigen Interrupt-Service-Routine kann Zeit vergehen. Dieses Intervall wird hierbei als Interrupt-Latenz bezeichnet. F{\"u}r Echtzeitsysteme ist vor allem relevant, welche Verz{\"o}gerung unter schlechten Bedingungen auftreten kann (Laplante \& Ovaska, 2012, S. 396--397).

Die Latenz wird sowohl durch die Hardware als auch durch den aktuellen Systemzustand beeinflusst. Bspw. kann eine softwareseitige Interrupt-Sperre die Behandlung verz{\"o}gern. Auch die Bedingungen f{\"u}r die Unterbrechung des gerade ausgef{\"u}hrten Maschinenbefehls spielen eine Rolle. Deshalb ist die Interrupt-Latenz nicht allein eine Eigenschaft der Anwendungsfunktion, die sp{\"a}ter in der ISR ausgef{\"u}hrt wird (Laplante \& Ovaska, 2012, S. 396--397).

Abschnitt 2.5.1c beschreibt die Interrupt-Infrastruktur von Zephyr. Wie die Interrupt-Anforderung auf der Zielplattform softwareseitig {\"u}ber den NVIC ausgel{\"o}st und ihre Latenz gemessen wird, wird erst in Kapitel 5.2 behandelt (Zephyr Project, 2024, \glqq{}Interrupts\grqq{}).

\subsection{Timer-Jitter}
\textbf{P{\"u}nktlichkeit \& Jitter }P{\"u}nktlichkeit beschreibt bei Laplante und Ovaska, wie nah die Antwortzeiten an einen vorgesehenen Antwortzeitpunkt liegen. F{\"u}r eine Anwendung k{\"o}nnen dabei sowohl zu fr{\"u}he als auch zu sp{\"a}te Reaktionen relevant sein. Abweichungen entstehen unter anderem durch die zeitlichen Eigenschaften der Hardware- und Softwarekomponenten (Laplante \& Ovaska, 2012, S. 8--9).

Wenn die Antwortzeiten zwischen wiederholten Vorg{\"a}ngen schwanken, dann wird diese zeitliche Variation Jitter genannt. F{\"u}r periodische Vorg{\"a}nge kann das Konzept auf die Abst{\"a}nde zwischen aufeinanderfolgenden Ausf{\"u}hrungen {\"u}bertragen werden. In dieser Arbeit ist die festgelegte Timerperiode der Bezugspunkt f{\"u}r die Untersuchung von solchen Intervallen. Die Messgr{\"o}{\ss}e und ihre Erfassung werden in Kapitel 5.3 erl{\"a}utert (Laplante \& Ovaska, 2012, S. 8--9).

Auf den periodischen Zephyr-Timer {\"u}bertragen bedeutet Jitter in dieser Arbeit die Abweichung des gemessenen Abstands zwischen zwei Callback-Messpunkten von der vorgesehenen Timerperiode. Die bestimmte Messgr{\"o}{\ss}e wird in Kapitel 5.3 festgelegt (Laplante \& Ovaska, 2012, S. 8--9; Zephyr Project, 2024, \glqq{}Timers\grqq{}).

\subsection{Heap-Allokation/ Fragmentierung}
\textbf{Dynamische Speicherverwaltung }Bei der dynamischen Speicherverwaltung wird Speicher w{\"a}hrend der Ausf{\"u}hrung angefordert und wieder freigegeben. In Echtzeitsystemen muss dabei sowohl die verf{\"u}gbare Kapazit{\"a}t ber{\"u}cksichtigt werden als auch der Aufwand der Verwaltung und das Verhalten bei fehlendem Speicher. Laplante und Ovaska verbinden die Speicherverwaltung deshalb mit der Vorhersagbarkeit des Gesamtsystems (2012, S. 127--132).

F{\"u}r diese Untersuchung sind zwei Aspekte relevant, darunter die Dauer von einzelnen Allokations- und Freigabeoperationen und die Verteilung des verf{\"u}gbaren Speichers nach wiederholten Anforderungen und Freigaben. Diese werden in getrennten Benchmarks untersucht.

\textbf{Fragmentierung }Wenn Speicher mehrmals zugewiesen und wieder freigegeben wird, kann der verf{\"u}gbare Speicherbereich aus belegten und freien Abschnitten bestehen. Wenn es insgesamt genug freien Speicher gibt, eine Anfrage jedoch nicht bearbeitet werden kann, weil es keinen zusammenh{\"a}ngenden Block der angeforderten Gr{\"o}{\ss}e gibt, dann nennt man dies externe Fragmentierung. Auf der anderen Seite kann eine Anfrage einen gr{\"o}{\ss}eren, fest vorgegebenen Speicherblock zugewiesen bekommen, als gebraucht wird. Der {\"u}brigbleibende Teil wird dabei nicht benutzt. Dies wird interne Fragmentierung genannt. Dabei kann eine Kompaktierung des Speichers eine externe Fragmentierung beheben. Dies ist jedoch selbst rechenintensiv und ist daher in harten Echtzeitsystemen unpraktisch. (Laplante \& Ovaska, 2012, S. 131-132)

In Abschnitt 2.5.1e ist die Speicherverwaltung von Zephyr {\"u}ber \lstinline[style=fileinline]|k_heap| bzw. dem System-Heap durch \lstinline[style=fileinline]|k_malloc()| und \lstinline[style=fileinline]|k_free()| erkl{\"a}rt. Dabei handelt es sich um eine Umsetzung dieser Konzepte.

\subsection{RAM/ ROM}
\textbf{Speicherklassen }Halbleiterspeicher kann in zwei Klassen eingeteilt werden. Dabei handelt es sich um fl{\"u}chtigen Lese-Schreib-Speicher (Random Access Memory, RAM) und nichtfl{\"u}chtigen Speicher (Read-Only Memory, ROM). Klassischer ROM l{\"a}sst sich dabei nur bei der Herstellung beschreiben. Auf der anderen Seite gibt es Vertreter wie EEPROM und Flash, welche auch im eingebauten Zustand neu beschreiben werden k{\"o}nnen. Dies passiert jedoch langsamer als bei RAM und auch nur f{\"u}r eine begrenzte Anzahl an Schreibzyklen. In eingebetteten Systemen wird ROM/ Flash daher oft zum Ablegen von Programmcode benutzt, w{\"a}hrend RAM f{\"u}r Stack und Laufzeitdaten reserviert bleibt. (Laplante \& Ovaska, 2012, S. 36-37)

\textbf{Speicher-Fu{\ss}abdruck }Der Speicherbedarf kann f{\"u}r verschiedene Bereiche getrennt betrachtet werden, bspw. f{\"u}r Programmcode, Daten und Stacks. Dabei wird zwischen der absoluten Belegung und ihrem Anteil an der verf{\"u}gbaren Kapazit{\"a}t unterschieden. Teile des Programmspeicherbedarfs k{\"o}nnen hierbei schon aus dem erzeugten Programmabbild bestimmt werden. Auf der anderen Seite h{\"a}ngt die wahre Stack-Auslastung noch vom Ausf{\"u}hrungsverlauf ab, bspw. von verschachtelten Funktionsaufrufen und Interrupts (Laplante \& Ovaska, 2012, S. 408--409).

Daher reicht f{\"u}r die Untersuchung der verf{\"u}gbaren Ressourcen nicht nur die Gesamtsumme. Ein einzelner Speicherbereich kann seine Kapazit{\"a}t n{\"a}mlich erreichen, obwohl an anderer Stelle noch Speicher frei ist. In dieser Arbeit wird der statische RAM-/ROM-Fu{\ss}abdruck anhand der erzeugten Firmware-Abbilder untersucht. Diese Auswertung soll dabei von einer Messung der echten Stack- oder Heap-Auslastung w{\"a}hrend des Betriebs unterschieden werden (Laplante \& Ovaska, 2012, S. 408--409).

In Abschnitt 2.5.2d sind die Werkezeuge \lstinline[style=fileinline]|ram_report| und \lstinline[style=fileinline]|rom_report| vom Zephyr-Build-System erw{\"a}hnt worden. Diese benutzen genau diese Art der Speicher-Fu{\ss}abdruck-Analyse f{\"u}r fertig gebaute Firmware-Abbildungen um.

\section{Netzwerk- und Sicherheitskonzepte}
Die untersuchte IoT-Stack-Konfiguration beinhaltet OpenThread, 6LoWPAN, CoAP und DTLS mit Pre-Shared Keys. MCUmgr und MCUboot geh{\"o}ren auf der anderen Seite nur zur getrennten Demonstrationsanwendung \lstinline[style=fileinline]|iot_sensor|. Die daf{\"u}r relevanten Netzwerk-, Sicherheits- und Updatekonzepte werden in diesem Abschnitt erl{\"a}utert.

\subsection{Netzwerk-Stack und CoAP}
Das Constrained Application Protocol (CoAP) ist ein eigenes, f{\"u}r ressourcenbeschr{\"a}nkte Ger{\"a}te gemachtes Web-{\"U}bertragungsprotokoll. Dabei folgt es demselben Anfrage-Antwort-Modell wie HTTP, darunter auch Methoden wie GET, POST, PUT und DELETE. CoAP benutzt UDP und erweitert dies mit eigenen Nachrichtentypen und optionale Zuverl{\"a}ssigkeitsmechanismen. Confirmable-Nachrichten werden best{\"a}tigt und bei ausbleibender Best{\"a}tigung neu {\"u}bertragen. Non-confirmable-Nachrichten verlangen hierbei keine solche Best{\"a}tigung (Shelby et al., 2014, Abschnitte 2 und 4). Zudem kann CoAP {\"u}ber eine gesicherte Variante, dem URI-Schema \glqq{}coaps\grqq{}, laufen, bei welcher Nachrichten durch DTLS abgesichert werden (Shelby et al., 2014) (siehe 2.7.2). Bei der von Silva et al. (2019, S. 10376) beschriebenen Netzwerk-Stack-Architektur bildet CoAP dabei die Anwendungsschicht, w{\"a}hrend UDP, IPv6 und die 6LoWPAN-Adaptionsschicht (siehe 2.7.3) die Schichten darunter {\"u}bernehmen.

\subsection{Transportverschl{\"u}sselung durch TLS/ DTLS und Pre-Shared Keys}
Transport Layer Security (TLS) ist ein verbreitetes Protokoll zur Absicherung von Netzwerkverbindungen. Es braucht jedoch einen zuverl{\"a}ssigen, verbindungsorientierten Transportkanal wie TCP. Da CoAP auf dem verbindungslosen UDP basiert, wird daher Datagram Transport Layer Security (DTLS) benutzt. Dabei handelt es sich um eine Variante, die bewusst so nah wie m{\"o}glich an TLS gehalten ist, daf{\"u}r aber {\"u}ber weitere Mechanismen f{\"u}r den Umgang mit verlorenen oder vertauschten Paketen verf{\"u}gt (Rescorla und Modadugu, 2012). DTLS nutzt dabei f{\"u}r die Absicherung Kryptographie-Verfahren, die durch sogenannte Cipher Suites festgelegt werden. F{\"u}r ressourcenbeschr{\"a}nkte Ger{\"a}te ist dabei eine Absicherung {\"u}ber einen davor ausgetauschten, gemeinsamen Schl{\"u}ssel (Pre-Shared Key, PSK) verbreitet, da daf{\"u}r kein aufwendiger Zertifikats- oder Schl{\"u}sselaustausch zur Laufzeit gebraucht wird. Die Cipher Suite \lstinline[style=fileinline]|TLS_PSK_WITH_AES_128_CCM_8| ist dabei die PSK-Absicherungsvariante, die f{\"u}r Constrained-Ger{\"a}te unterst{\"u}tzt werden muss (Shelby et al., 2014). Dabei gibt es sowohl einen Nach- als auch einen Vorteil. Die in dieser Arbeit benutzte Cipher Suite \lstinline[style=fileinline]|TLS_PSK_WITH_AES_128_CCM_8| bietet keine Perfect Forward Secrecy, besitzt daf{\"u}r aber einen vergleichsweise kleinen Speicherbedarf und eine f{\"u}r ressourcenbeschr{\"a}nkte Ger{\"a}te passende Ausf{\"u}hrungseffizienz (Tschofenig \& Fossati, 2016, Abschn. 9). Lee et al. implementierten f{\"u}r Zephyr unter anderem \lstinline[style=fileinline]|PSK_WITH_AES128_CCM8| und \lstinline[style=fileinline]|ECDHE_ECDSA_WITH_AES128_CCM8|. Ihre Hardwaretests benutzen ein FRDM-K64F mit Cortex-M4, 1 MiB Flash und 256 KiB RAM. Trotz gleicher Speichergr{\"o}{\ss}en ist dieses Board wegen weiterer Hardware- und Softwareunterschiede keine direkt vergleichbare Messplattform (Lee et al., 2018, S. 1292--1294).

\subsection{Thread, 6LoWPAN und IEEE 802.15.4}
IEEE 802.15.4 ist ein Standard f{\"u}r stromsparende Funknetze mit niedriger Datenrate. Dieser Standard erlaubt dabei nur kleine Nutzlasten pro {\"U}bertragungsrahmen. Diese sind kleiner als die von IPv6 mindestens geforderte Paketgr{\"o}{\ss}e. Daher wird eine weitere Zwischenschicht gebraucht, die IPv6-Pakete auf mehrere kleine 802.15.4-Rahmen aufteilt und dabei die IPv6-Header komprimiert, da sie sonst zu viel Platz verbrauchen w{\"u}rden. Diese Zwischenschicht wird 6LoWPAN genannt (Montenegro et al., 2007). Thread ist hierbei ein darauf aufbauendes, IPv6-basiertes Mesh-Netzwerkprotokoll, das genau diese 6LoWPAN-Kompression {\"u}ber IEEE-802.15.4.Funkverbindung benutzt, um ein IPv6-basiertes Mesh-Netz aus ressourcenbeschr{\"a}nkten Ger{\"a}ten zu bilden (Zephyr Project, 2024, \glqq{}Thread protocol\grqq{}). Damit aber ein solches Thread-Netz mit dem klassischen Internet kommunizieren kann, wird ein weiterer sogenannter Border Router gebraucht, welcher zwischen dem Thread-Mesh und einem externen IP-Netz vermittelt (Zephyr Project, 2024, \glqq{}Thread protocol\grqq{}). Die f{\"u}r diese Arbeit ausgesuchte IoT-Stack-Konfiguration benutzt dabei das nRF52840 DK als sogenanntes Minimal Thread Device (MTD), also als ressourcenschonenden Netzknoten ohne eigene Routing-Funktion innerhalb vom Mesh (siehe 4.4.3).

\subsection{Firmware-Updates: MCUmgr und MCUboot}
MCUboot verwaltet Firmware-Slots und kann dabei entsprechend konfigurierte und signierte Firmware-Abbilder vor dem Start pr{\"u}fen. MCUmgr stellt Managementfunktionen bereit, {\"u}ber die neue Images {\"u}bertragen und verwaltet werden k{\"o}nnen. Zephyr unterst{\"u}tzt daf{\"u}r getrennte SMP-Transporte {\"u}ber Bluetooth LE, serielle Verbindungen und UDP. In der Demonstrationsanwendung dieser Arbeit wird UDP {\"u}ber den vorhandenen IPv6-/ Thread-Stack verwendet (Zephyr Project, 2024, \glqq{}MCUmgr\grqq{} und \glqq{}Device Firmware Upgrade\grqq{}).
